\documentclass[12pt]{article}
\usepackage[utf8]{inputenc}
\usepackage[T1]{fontenc}
\usepackage[margin=1in]{geometry}
\usepackage{setspace}
\usepackage{xcolor}
\usepackage[colorlinks=true, urlcolor=blue!60!black, linkcolor=blue!60!black]{hyperref}
\setlength{\parindent}{2em}
\setlength{\parskip}{1.25em}
\renewcommand{\baselinestretch}{1.0}
\begin{document}
\begin{center}
{\LARGE \textit{Rucho v. Common Cause}, 588 U.S. 684 (2019)\par}
\vspace{0.8em}
{\large Case Brief\par}
\end{center}
\vspace{1em}

\section{Memory Jogger}

Partisan gerrymandering claims present nonjusticiable political questions beyond the jurisdiction of the federal courts because the Constitution supplies no clear, manageable, and politically neutral legal standard for deciding how much partisanship in districting is too much (\href{https://www.courtlistener.com/opinion/4633469/rucho-v-common-cause/}{CourtListener}; \href{https://www.law.cornell.edu/supremecourt/text/18-422}{Cornell LII}).

\section{Detailed Case Facts}

\textit{Rucho v. Common Cause} consolidated challenges to two overtly partisan congressional redistricting maps---one drawn by Republicans in North Carolina and the other by Democrats in Maryland (\href{https://www.law.cornell.edu/supremecourt/text/18-422}{Cornell LII}; \href{https://www.supremecourt.gov/opinions/18pdf/18-422_9ol1.pdf}{Supreme Court Slip Op.}).

In North Carolina (\textit{Rucho v. Common Cause}, No. 18--422), after a federal district court struck down two districts in the State's post-2010 congressional map as unconstitutional racial gerrymanders, the Republican-controlled General Assembly enacted a remedial congressional redistricting plan in 2016. The Republican co-chairs of the redistricting committee, State Senator Robert Rucho and State Representative David Lewis, instructed Republican mapmaker Dr. Thomas Hofeller to use precinct-level political data to draw a map that would maintain a 10--3 Republican advantage in North Carolina's 13-seat U.S. House delegation. The committee formally adopted ``Partisan Advantage'' as a redistricting criterion, and Representative Lewis openly explained that he drew the map for 10 Republicans and 3 Democrats only because he did not believe it was possible to draw an 11--2 map, adding that he thought electing Republicans was better for the country. Adopted on a party-line vote, the 2016 Plan delivered 10 Republican seats and 3 Democratic seats in 2016 (with Republicans winning 53\% of the statewide vote) and 9 Republican seats and 3 Democratic seats in 2018 (with Republicans winning 50\% of the vote, and the remaining seat ordered to a new election following fraud allegations) (\href{https://www.law.cornell.edu/supremecourt/text/18-422}{Cornell LII}; \href{https://harvardlawreview.org/print/vol-133/rucho-v-common-cause/}{Harvard Law Review}).

In Maryland (\textit{Lamone v. Benisek}, No. 18--726), the Democratic-controlled legislature and Democratic Governor Martin O'Malley redrew Maryland's eight congressional districts following the 2010 census. Governor O'Malley appointed an advisory committee and enlisted Congressman Steny Hoyer---a self-described ``serial gerrymanderer''---to direct political analyst Eric Hawkins to craft a map that would protect Democratic incumbents and flip the Sixth Congressional District, which a Republican had held for nearly two decades, to create a 7--1 Democratic delegation. Although equal-population requirements necessitated removing only about 10,000 residents from the Sixth District, the 2011 Plan removed roughly 360,000 voters and added 350,000 new voters, reducing registered Republicans in the district by about 66,000 and increasing registered Democrats by about 24,000. Enacted on a party-line vote, the map flipped the Sixth District in 2012, and Democrats won 7 of Maryland's 8 seats in every election from 2012 through 2018 while never exceeding 65\% of the statewide congressional vote (\href{https://www.law.cornell.edu/supremecourt/text/18-422}{Cornell LII}).

\section{Procedural History}

\textbf{North Carolina Litigation.} In August 2016, the North Carolina Democratic Party, Common Cause, and 14 individual voters sued legislative leaders and state officials in the U.S. District Court for the Middle District of North Carolina; the League of Women Voters of North Carolina and a dozen voters filed a parallel suit that was consolidated. Plaintiffs alleged that the 2016 Plan violated the Fourteenth Amendment's Equal Protection Clause, the First Amendment, Article I, \S~2, and the Elections Clause (Article I, \S~4, cl. 1). After a four-day trial, a three-judge panel unanimously struck down the plan under Equal Protection and Article I, and held 2--1 (Judge Osteen dissenting) that it violated the First Amendment (\textit{Common Cause v. Rucho}, 279 F. Supp. 3d 587 (M.D.N.C. 2018)). On direct appeal under 28 U.S.C. \S~1253, the Supreme Court vacated and remanded for reconsideration of standing in light of \textit{Gill v. Whitford}, 585 U.S. 48 (2018). On remand, the three-judge panel found standing, invalidated 12 of the 13 districts under Equal Protection, reaffirmed its First Amendment (again over Judge Osteen's dissent) and Article I rulings, and enjoined use of the map after the November 2018 election (\textit{Common Cause v. Rucho}, 318 F. Supp. 3d 777 (M.D.N.C. 2018)) (\href{https://www.law.cornell.edu/supremecourt/text/18-422}{Cornell LII}; \href{https://clearinghouse.net/case/43526/}{Civil Rights Litigation Clearinghouse}).

\textbf{Maryland Litigation.} In November 2013, Maryland Republican voters sued state election officials in the U.S. District Court for the District of Maryland, challenging the 2011 Plan under the First Amendment, Article I, \S~2, and the Elections Clause. After earlier proceedings on preliminary injunctive relief (\textit{Benisek v. Lamone}, 585 U.S. 155 (2018) (per curiam)), a three-judge panel granted summary judgment to the plaintiffs on First Amendment grounds---finding burdens on both representational and associational rights---and permanently enjoined the 2011 Plan (\textit{Benisek v. Lamone}, 348 F. Supp. 3d 493 (D. Md. 2018)) (\href{https://www.law.cornell.edu/supremecourt/text/18-422}{Cornell LII}; \href{https://harvardlawreview.org/print/vol-133/rucho-v-common-cause/}{Harvard Law Review}).

\textbf{Supreme Court Review.} Both sets of state defendants appealed directly to the Supreme Court under 28 U.S.C. \S~1253. The Court postponed jurisdiction, consolidated the appeals, heard oral argument on March 26, 2019, and issued its decision on June 27, 2019, vacating both judgments and remanding with instructions to dismiss for lack of jurisdiction (\href{https://www.law.cornell.edu/supremecourt/text/18-422}{Cornell LII}; \href{https://www.supremecourt.gov/opinions/18pdf/18-422_9ol1.pdf}{Supreme Court Slip Op.}).

\section{Judicial Votes}

The Supreme Court decided the consolidated cases by a \textbf{5--4} vote, with no concurring opinions (\href{https://www.courtlistener.com/opinion/4633469/rucho-v-common-cause/}{CourtListener}; \href{https://www.law.cornell.edu/supremecourt/text/18-422}{Cornell LII}):

\begin{itemize}
\item \textbf{Majority (5):} Chief Justice John G. Roberts, Jr. (author), joined by Justices Clarence Thomas, Samuel A. Alito, Jr., Neil M. Gorsuch, and Brett M. Kavanaugh.
\item \textbf{Concurrences (0):} None.
\item \textbf{Dissent (4):} Justice Elena Kagan (author), joined by Justices Ruth Bader Ginsburg, Stephen G. Breyer, and Sonia Sotomayor.
\end{itemize}

\section{Holding}

\textbf{Judgment.} The judgments of the U.S. District Court for the Middle District of North Carolina and the U.S. District Court for the District of Maryland were vacated, and both cases were remanded with instructions to dismiss for lack of jurisdiction (\href{https://www.law.cornell.edu/supremecourt/text/18-422}{Cornell LII}).

\textbf{Rule Announced.} Claims of unconstitutional partisan gerrymandering---whether framed under the Equal Protection Clause, the First Amendment, Article I, \S~2, or the Elections Clause---present nonjusticiable political questions beyond the reach of the federal courts under Article III because there are no judicially discoverable and manageable standards for deciding when permissible consideration of partisan advantage in districting becomes unconstitutional (\href{https://www.law.cornell.edu/supremecourt/text/18-422}{Cornell LII}).

\textbf{Holding vs. Dicta.} The core holding is strictly jurisdictional and federal: Article III federal courts lack jurisdiction to adjudicate partisan gerrymandering claims. In contrast, the Court's observations that excessive partisan gerrymandering is ``incompatible with democratic principles'' and ``leads to results that reasonably seem unjust,'' as well as its survey of state constitutional amendments, independent redistricting commissions, state-court rulings (such as Florida's \textit{League of Women Voters v. Detzner}), and pending congressional bills under the Elections Clause, are non-binding dicta explaining why nonjusticiability does not condone the practice or leave complaints without non-federal avenues of reform. Likewise, the Court's suggestion that claims invoking the core principle that voters should choose their representatives sound more naturally under the nonjusticiable Guarantee Clause (Article IV, \S~4) reinforces, but is secondary to, its lack-of-manageable-standards holding (\href{https://www.law.cornell.edu/supremecourt/text/18-422}{Cornell LII}; \href{https://harvardlawreview.org/print/vol-133/rucho-v-common-cause/}{Harvard Law Review}).

\section{Analysis of Opinions}

\subsection{Majority Opinion (Chief Justice Roberts)}

Chief Justice Roberts grounded the Court's analysis in Article III's ``Cases'' and ``Controversies'' limitation and the political question doctrine articulated in \textit{Baker v. Carr}, 369 U.S. 186 (1962), specifically the absence of ``judicially discoverable and manageable standards.'' He began with Founding-era history, noting that partisan gerrymandering predated Independence (including Patrick Henry's 1789 attempt to gerrymander James Madison out of Congress and Elbridge Gerry's 1812 Massachusetts map) and that the Framers addressed electoral manipulation through the Elections Clause (Article I, \S~4, cl. 1), assigning districting primarily to state legislatures checked by Congress rather than federal courts. Although the Court rejected the appellants' argument that the Elections Clause completely bars judicial review of all districting issues---preserving federal judicial review of one-person, one-vote and racial gerrymandering claims---it emphasized that those domains are qualitatively different: population equality is administrable by arithmetic, and racial gerrymandering demands the elimination of an inherently suspect racial classification, whereas precedent recognizes that ``a jurisdiction may engage in constitutional political gerrymandering'' (\href{https://www.law.cornell.edu/supremecourt/text/18-422}{Cornell LII}).

Because some partisan consideration is constitutionally permissible, the central question is one of degree: ``How much is too much?'' Roberts argued that partisan gerrymandering claims inevitably rest on an unenumerated desire for proportional representation, which the Constitution does not require. Even defining ``fairness''---whether by maximizing competitive seats, creating proportional safe seats, or adhering to traditional criteria that reflect ``natural political geography''---requires an inherently political, not legal, judgment. The majority then systematically rejected every proposed test: (1) the North Carolina Equal Protection test (predominant intent, durable dilution, lack of neutral justification) failed because permissible partisan intent does not become unlawful simply by predominating, and predicting future electoral durability asks judges to ``beat the point spread'' despite historical reversals after \textit{Davis v. Bandemer} and \textit{Vieth v. Jubelirer}; (2) the First Amendment retaliation/association test failed because districting plans do not restrict speech or association, and treating any partisan disadvantage as viewpoint discrimination would outlaw all political considerations; (3) the dissent's ``extreme outlier'' baseline derived from a State's own nonpartisan criteria was deemed arbitrary because federal constitutional meaning would shift with state-chosen criteria and still left unanswered how far from the median map is ``too much''; and (4) Article I, \S~2 and the Elections Clause provide no judicially enforceable limit on partisan considerations. Concluding that federal judges have no commission to reallocate political power, Roberts pointed to state constitutional reforms, independent commissions, state courts, and Congress's Elections Clause power as the proper avenues for reform (\href{https://www.law.cornell.edu/supremecourt/text/18-422}{Cornell LII}; \href{https://www.supremecourt.gov/opinions/18pdf/18-422_9ol1.pdf}{Supreme Court Slip Op.}).

\subsection{Dissenting Opinion (Justice Kagan)}

Justice Kagan, joined by Justices Ginsburg, Breyer, and Sotomayor, opened by characterizing the ruling as the first time in American history that the Court refused to remedy an acknowledged constitutional violation on the ground that it exceeded judicial capabilities. Recounting the explicit partisan intent and dramatic electoral skews in both North Carolina and Maryland, Kagan argued that extreme partisan gerrymandering turns popular sovereignty upside down---allowing politicians to choose their voters rather than voters choosing their representatives---and violates both the Equal Protection Clause (by intentionally diluting disfavored voters' ballots through packing and cracking) and the First Amendment (by penalizing citizens for their political views and chilling associational effectiveness). She stressed that modern big data, precinct-level analytics, and computing power make contemporary gerrymanders far more durable and destructive than Founding-era linedrawing (\href{https://www.law.cornell.edu/supremecourt/text/18-422}{Cornell LII}).

Kagan contended that lower federal courts had already converged on a neutral, manageable, three-part framework---predominant intent to entrench, substantial vote-dilutive effect, and lack of a legitimate nonpartisan justification---that did not require judges to impose proportional representation or personal ideals of fairness. Instead, by using computer-generated ensembles of thousands of simulated maps (such as 3,000 or 24,518 maps in North Carolina) built from each State's own physical geography, political geography, and nonpartisan redistricting criteria, courts could objectively identify ``extreme outliers'' without second-guessing ordinary politics. As for the majority's ``how much is too much?'' objection, Kagan answered ``This much is too much,'' noting that courts routinely adjudicate matters of degree (such as ``substantial'' harm or anticompetitive effects) without rigid numerical thresholds. Finally, she dismissed the majority's reliance on Congress, ballot initiatives, and state courts as unrealistic---politicians who benefit from gerrymandering will not voluntarily dismantle it, fewer than half the States permit voter initiatives, and if state courts can apply broad state constitutional clauses to strike down gerrymanders, federal courts can do so as well (\href{https://www.law.cornell.edu/supremecourt/text/18-422}{Cornell LII}).

\section{Examples: Future Applications}

\begin{enumerate}
\item \textbf{Same-Side (Nonjusticiable in Federal Court): State Legislative Partisan Gerrymander.} A state legislature uses machine-learning software to draw state senate districts that lock in a 65\% supermajority for the ruling party even when it wins only 48\% of the statewide vote, and voters sue in federal district court alleging Fourteenth Amendment vote dilution and First Amendment retaliation. Under \textit{Rucho}, the federal court must dismiss for lack of jurisdiction because partisan gerrymandering claims---whether involving congressional or state legislative maps---present nonjusticiable political questions.
\item \textbf{Same-Side (Nonjusticiable in Federal Court): Mid-Decade Partisan Redistricting.} Midway through a decade, a newly elected state legislative majority redraws congressional districts solely to oust two opposition incumbents, while maintaining exact population equality and avoiding racial classifications. Voters sue in federal court under Article I, \S~2 and the Elections Clause. Under \textit{Rucho}, this claim is barred as a nonjusticiable political question because neither Article I nor the Fourteenth Amendment supplies a manageable federal standard to police partisan motivation in districting.
\item \textbf{Opposite-Side (Justiciable): State-Court Challenge Under a State Constitution.} Voters challenge a state legislature's aggressively gerrymandered congressional map in state court under a state constitutional ``Free and Equal Elections'' clause or an explicit state constitutional ban on partisan favoritism. \textit{Rucho} explicitly leaves state-court enforcement of state constitutional and statutory redistricting limits intact, and the Supreme Court later confirmed in \href{https://www.supremecourt.gov/opinions/22pdf/21-1271_3f14.pdf}{\textit{Moore v. Harper}, 600 U.S. 1 (2023)}, that the federal Elections Clause does not insulate state legislatures from ordinary state judicial review under state constitutions.
\item \textbf{Opposite-Side (Justiciable): Racial Gerrymandering and One-Person, One-Vote Violations.} A state legislature draws a congressional map with a 12\% population deviation between districts, or intentionally uses race as the predominant factor to pack Black voters into a single district without a compelling justification under the Voting Rights Act. Under \textit{Rucho}, federal courts retain full jurisdiction to adjudicate both one-person, one-vote claims (administrable mathematically) and racial gerrymandering claims (targeting an inherently suspect racial classification), because neither requires courts to measure how much partisanship is ``too much.''
\item \textbf{Fence-Sitter: Conflated Race-and-Party Gerrymandering Claims.} In a State where racial demographics and partisan affiliation are tightly correlated, a legislature moves tens of thousands of Black Democratic voters out of a competitive congressional district to make the seat safer for Republicans, claiming purely partisan motives. Because racial gerrymandering remains justiciable while partisan gerrymandering does not, courts face a razor-thin boundary: plaintiffs must disentangle race from party and overcome a presumption of legislative good faith, a tension later highlighted in \href{https://www.acslaw.org/scotus_update/alexander-v-south-carolina-state-conference-of-the-naacp/}{\textit{Alexander v. South Carolina State Conference of the NAACP}, 602 U.S. 1 (2024)}, where the Court demanded strong proof (including an alternative map) that race rather than \textit{Rucho}-protected partisanship drove the lines.
\end{enumerate}

\section{Critique}

\textbf{Originalist and Textualist Perspectives.} Many conservative and originalist commentators praised \textit{Rucho}'s outcome as faithful to the Constitution's text and separation of powers, emphasizing that the Elections Clause expressly vests oversight of state districting abuses in Congress, not Article III courts (\href{https://fedsoc.org/case/rucho-v-common-cause}{Federalist Society}). Notably, Professor John O. McGinnis characterized \textit{Rucho} as an exercise in ``silent originalism'': while agreeing with the result, he argued that the majority did not even need to rely on the political question doctrine because, as an original matter, the Fourteenth Amendment's Equal Protection Clause did not create a substantive right against partisan vote dilution (nor even the one-person, one-vote rule of \textit{Reynolds v. Sims}, which Roberts conspicuously avoided citing) (\href{https://lawliberty.org/silent-originalism-and-the-reweighting-of-precedent/}{Law \& Liberty}). Other formalists, however, question whether labeling a claim a nonjusticiable ``political question'' under Article III is doctrinally coherent when the real issue is whether the Constitution confers a substantive right at all.

\textbf{Progressive and Doctrinal Perspectives.} Progressive and election-law scholars sharply criticized \textit{Rucho} as an abdication of the judiciary's representation-reinforcing duty in precisely the domain where political self-correction fails because entrenched incumbents control the rules of electoral competition. The \href{https://harvardlawreview.org/print/vol-133/rucho-v-common-cause/}{Harvard Law Review} noted that although Chief Justice Roberts framed \textit{Rucho} in ``classical'' Article III terms, the decision actually rested heavily on ``prudential'' concerns---fear of delegitimizing remedial mapdrawing and optimism about state-level reforms---and created a doctrinal anomaly in which state courts might remain free to adjudicate federal claims that federal courts cannot hear. Scholars also challenged the majority's insistence that computer-generated ensemble methods lack manageability, arguing that \textit{Rucho} effectively greenlights technological partisan entrenchment and incentivizes litigants to repackage partisan disputes as racial gerrymandering claims (\href{https://harvardlawreview.org/print/vol-133/rucho-v-common-cause/}{Harvard Law Review}).

\textbf{Underlying Competing Values.} At bottom, \textit{Rucho} reflects an enduring clash between \textit{judicial restraint and institutional legitimacy} (preventing unelected federal judges from becoming perpetual referees of partisan power allocation without clear textual rules) and \textit{democratic self-government and individual voting equality} (preventing self-interested legislators from insulating themselves from electoral accountability).

\section{Key Quotations}

\begin{quote}
``Excessive partisanship in districting leads to results that reasonably seem unjust. But the fact that such gerrymandering is `incompatible with democratic principles,' \textit{Arizona State Legislature}, 576 U.\,S., at 791 (brackets omitted), does not mean that the solution lies with the Federal Judiciary. We conclude that partisan gerrymandering claims present political questions beyond the reach of the federal courts.''
\end{quote}
\noindent --- Chief Justice Roberts, writing for the Court, 588 U.S.\ at 718 (\href{https://www.law.cornell.edu/supremecourt/text/18-422}{Cornell LII}; \href{https://www.courtlistener.com/opinion/4633469/rucho-v-common-cause/}{CourtListener}).

\begin{quote}
``Our conclusion does not condone excessive partisan gerrymandering. Nor does our conclusion condemn complaints about districting to echo into a void.''
\end{quote}
\noindent --- Chief Justice Roberts, writing for the Court, 588 U.S.\ at 719 (\href{https://www.law.cornell.edu/supremecourt/text/18-422}{Cornell LII}).

\begin{quote}
``No one can accuse this Court of having a crabbed view of the reach of its competence. But we have no commission to allocate political power and influence in the absence of a constitutional directive or legal standards to guide us in the exercise of such authority. `It is emphatically the province and duty of the judicial department to say what the law is.' \textit{Marbury v. Madison}, 1 Cranch, at 177. In this rare circumstance, that means our duty is to say `this is not law.' ''
\end{quote}
\noindent --- Chief Justice Roberts, writing for the Court, 588 U.S.\ at 721 (\href{https://www.law.cornell.edu/supremecourt/text/18-422}{Cornell LII}; \href{https://www.supremecourt.gov/opinions/18pdf/18-422_9ol1.pdf}{Supreme Court Slip Op.}).

\begin{quote}
``For the first time ever, this Court refuses to remedy a constitutional violation because it thinks the task beyond judicial capabilities.''
\end{quote}
\noindent --- Justice Kagan, dissenting, 588 U.S.\ at 721 (\href{https://www.law.cornell.edu/supremecourt/text/18-422}{Cornell LII}).

\begin{quote}
``Of all times to abandon the Court's duty to declare the law, this was not the one. The practices challenged in these cases imperil our system of government. Part of the Court's role in that system is to defend its foundations. None is more important than free and fair elections. With respect but deep sadness, I dissent.''
\end{quote}
\noindent --- Justice Kagan, dissenting (\href{https://www.law.cornell.edu/supremecourt/text/18-422}{Cornell LII}; \href{https://www.courtlistener.com/opinion/4633469/rucho-v-common-cause/}{CourtListener}).

\end{document}
